\chapter{Конструкторский раздел}

В данном разделе представляются требования к серверу, архитектура сервера, алгоритмы работы сервера, выбор типов и структур данных и классы эквивалентности для тестирования.

\section{Требования к серверу}

Разрабатываемый сервер должен отвечать следующим требованиям:
\begin{enumerate}
	\item поддержка запросов \texttt{GET} и \texttt{HEAD};
	\item поддержка статусов, соответствующих успешному выполнению и ошибкам на стороне клиента;
	\item поддержка передачи файлов до 128 Мбайт;
	\item использование механизма мультиплексирования;
	\item использование логирования.
\end{enumerate}

\section{Архитектура сервера}

Архитектура сервера представлена в нотации моделирования архитектуры программных систем C4~\cite{c4}.  

\subsection{Диаграмма контекста}

Диаграмма контекста демонстрирует, как программная система взаимодействует с внешними сущностями. 

Диаграмма контекста разрабатываемого сервера представлена на рисунке~\ref{img:context}.

\includeimage
{context}
{f} % Обтекание (без обтекания)
{h}
{\textwidth} % Ширина рисунка
{Диаграмма контекста}

\subsection{Диаграмма контейнеров}

Диаграмма контейнеров описывает верхнеуровневую архитектуру и технологии. Используется для понимания технологического стека и разделения зон ответственности.

Диаграмма контейнеров разрабатываемого сервера представлена на рисунке~\ref{img:container}.

\includeimage
{container}
{f} % Обтекание (без обтекания)
{h}
{\textwidth} % Ширина рисунка
{Диаграмма контейнеров}

\section{Алгоритмы работы сервера}

\subsection{Алгоритм запуска сервера}

\subsection{Алгоритм прослушивания запросов к серверу}

\subsection{Алгоритм обработки запроса к серверу}



\section{Выбор типов и структур данных}

\section{Классы эквивалентности для тестирования}

\section*{Вывод}